\section{Conclusion}

Angle-resolved Fourier-operator metasurface design is better posed as a response-space inverse problem than as a sequence of disconnected task optimizations. That is the central conclusion of this manuscript. By treating the target as a coupled field over angle, wavelength, and polarization, the study reframes second-order differentiation, higher-order differentiation, angular filtering, and polarization-dependent operation as different instances of one physical specification language. The associated inverse-design strategy is naturally generative because the feasible set is multimodal, and it is necessarily physics-consistent because surrogate agreement alone does not certify operator behavior.

Within that formulation, the present results support a design loop built from simulation-grounded supervision, surrogate-assisted candidate concentration, conditional diffusion-based inverse generation, and RCWA-based verification and reranking. The scientific contribution is therefore not a catalog of separate devices. It is a response-centered inverse-design paradigm for freeform Fourier-operator metasurfaces in which target specification, candidate generation, and final selection remain tied to the verified scattering behavior of the device. From that perspective, future progress is likely to depend less on accumulating isolated operator demonstrations and more on learning, sampling, and verifying structured optical response manifolds.
